\section{The actual rounded canonical module}

\begin{lemma}[An actual line atlas gives native invertibility]
\label{lem:canonical-line-atlas-native-invertible}
\lean{EnvelopeProperty.CanonicalGeometry.LocalBundleFrame.pullbackIso,EnvelopeProperty.CanonicalGeometry.LineBundleAtlas.isInvertible}
\leanok
An actual open line atlas of an original module sheaf gives a covering by
actual pullback frames, hence invertibility of that same module sheaf.
\end{lemma}
\begin{proof}
For each atlas open, the equivalence from its over-site to the actual open
subscheme carries the given whole frame to a restriction frame. Its unit
comparison identifies the transported structure module with the original
structure module of the open subscheme. The actual restriction--pullback
isomorphism turns it into a pullback frame. The atlas's pointwise covering
property says exactly that these opens have supremum the whole scheme.
\end{proof}

\begin{theorem}[Native invertibility of the original canonical sheaf]
\label{thm:canonical-original-native-invertible}
\uses{lem:canonical-line-atlas-native-invertible}
\lean{EnvelopeProperty.CanonicalGeometry.canonicalLineSheaf_isInvertible}
\leanok
If $f:Y\to\operatorname{Spec}k$ is smooth of relative dimension $n$, its
original sheafified exterior cotangent sheaf $\bigwedge^n\Omega_{Y/k}$ is
an invertible module sheaf.
\end{theorem}
\begin{proof}
The actual standard-smooth affine charts of $f$ supply bases of the
original cotangent modules. Their top wedges give the actual line atlas
of the original sheafified exterior cotangent presheaf. Apply the preceding
lemma to this atlas. No alternative canonical object is chosen.
\end{proof}

\begin{lemma}[The original rounded ideal module]
\label{lem:rounded-snc-original-module}
\uses{thm:rounded-snc-original-ideal}
\lean{EnvelopeProperty.Char0.RoundedCanonicalModule.idealModule,EnvelopeProperty.Char0.RoundedCanonicalModule.idealModule_affineSectionsEquiv,EnvelopeProperty.Char0.RoundedCanonicalModule.idealModule_isInvertible}
\leanok
The module sheaf of the original rounded ideal is actually invertible.
On every original affine open its entire section module is linearly
isomorphic to the original rounded affine ideal; on adapted charts that
ideal is the original rounded coordinate monomial ideal.
\end{lemma}
\begin{proof}
Use the genuine module sheaf of the whole rounded ideal. Its section
comparison with the original affine ideal is the actual ideal-module
affine equivalence. The already proved rounded-ideal invertibility gives
regular local equations. Multiplication by each such equation is a
linear isomorphism onto that same original ideal, so the actual affine
module comparisons give local pullback frames. This proves invertibility
of the original ideal module. The adapted-chart formula is the whole
rounded ideal's formula through this exact section equivalence.
\end{proof}

\begin{lemma}[Actual homomorphisms as sections of a tensor-dual line]
\label{lem:original-invertible-hom-section}
\uses{thm:scheme-original-invertible-tensor-exact}
\lean{EnvelopeProperty.SchemeInvertibleHomSection.equiv,EnvelopeProperty.SchemeInvertibleHomSection.recover,EnvelopeProperty.SchemeInvertibleHomSection.section_recover}
\leanok
For an original invertible module $A$ and any original module $B$, there
is a genuine bijection
\[
 \operatorname{Hom}(A,B)\simeq
 \operatorname{Hom}(\mathcal O,B\otimes A^\vee).
\]
Its inverse recovers the entire original homomorphism.
\end{lemma}
\begin{proof}
Tensoring on the right by the actual $A$ is an equivalence with inverse
tensoring on the right by the actual sheaf dual $A^\vee$. This equivalence
is built from the actual evaluation isomorphism and the tensor
associativity, unit and braiding comparisons. Its adjunction gives
$\operatorname{Hom}(\mathcal O\otimes A,B)\simeq
\operatorname{Hom}(\mathcal O,B\otimes A^\vee)$. Precompose the source
identification with the actual left unitor. The inverse laws of these two
bijections prove recovery on every whole original map.
\end{proof}

\begin{theorem}[The actual rounded relative canonical line]
\label{thm:rounded-relative-canonical-original-module}
\uses{thm:canonical-original-native-invertible,lem:rounded-snc-original-module}
\lean{EnvelopeProperty.Char0.RoundedCanonicalModule.absolute,EnvelopeProperty.Char0.RoundedCanonicalModule.relative,EnvelopeProperty.Char0.RoundedCanonicalModule.relativeRounded,EnvelopeProperty.Char0.RoundedCanonicalModule.absolute_isInvertible,EnvelopeProperty.Char0.RoundedCanonicalModule.relative_isInvertible,EnvelopeProperty.Char0.RoundedCanonicalModule.relativeRounded_isInvertible}
\leanok
Let $g:Y\to X$ be arbitrary, with both $X$ and $Y$ smooth of the same
relative dimension $n$ over $k$, and let $J$ be an actual SNC divisor ideal
on $Y$, and let $q\in\mathbb Q_{\ge0}$. The genuine module
\[
 \omega_{Y/k}\otimes(g^*\omega_{X/k})^\vee
       \otimes\mathcal O_Y(-\lfloor qF\rfloor),\qquad J=\mathcal O_Y(-F),
\]
is an invertible module sheaf on the WHOLE original $Y$.
\end{theorem}
\begin{proof}
The two original canonical sheaves are invertible by smoothness. Pullback
along the actual arbitrary $g$ preserves invertibility, as does the
genuine sheaf dual. Tensor the resulting actual line with the original
canonical line on $Y$ and the already constructed whole rounded ideal
module. Each operation preserves actual invertibility. This argument
requires no smoothness, flatness, properness or birationality of $g$.
\end{proof}

\begin{theorem}[The actual exterior differential as a relative canonical section]
\label{thm:canonical-original-differential-section}
\uses{thm:canonical-original-native-invertible,lem:original-invertible-hom-section}
\lean{EnvelopeProperty.Char0.RoundedCanonicalModule.differentialSection,EnvelopeProperty.Char0.RoundedCanonicalModule.differentialSection_recover}
\leanok
Let $f:X\to\operatorname{Spec}k$ be smooth of pure relative dimension
$n$, and let $g:Y\to X$ be any morphism, with $Y$ arbitrary. The actual
exterior differential
$g^*\bigwedge^n\Omega_{X/k}\to\bigwedge^n\Omega_{Y/k}$ defines a genuine
section of the original relative canonical module
$\omega_{Y/k}\otimes(g^*\omega_{X/k})^\vee$. This module is a line when
$Y$ is also smooth of the same pure relative dimension. The inverse
section--map comparison is exactly the entire original exterior differential map.
\end{theorem}
\begin{proof}
Apply the preceding Hom--section bijection to the actual arbitrary-map
exterior cotangent pullback morphism, using the original canonical sheaves
as its actual source and target. Invertibility of the source follows
from the original smooth canonical sheaf and arbitrary pullback. Recovery
is the inverse law of that genuine bijection; consequently no Jacobian,
section or image comparison is an extra premise here.
\end{proof}
